% =====================================================================
% BOZZA della sezione "Test case description".
% Questo file NON e' incluso nella tesi: quando la bozza va bene,
% il contenuto va al posto della sezione omonima in 02_metodologia_setup.tex
% =====================================================================

\section{Test case description}
\label{sec:val_testcase}

\subsection{Geometry and reference quantities}
\label{sec:val_geometry}

The geometry considered in this work is the wing/body/horizontal tail configuration of the CRM with the stabiliser set at \SI{0}{\degree} incidence (WBT0), in the full-scale version released for DPW-IV~\cite{vassberg2008}. The models tested in wind tunnels are scaled replicas of this geometry: the one used in the ETW and NTF campaigns has a mean aerodynamic chord of \SI{7.45}{in}, which corresponds to a scale factor of about 2.7\%~\cite{rivers2015}. Working on the full-scale geometry does not affect the comparison with the experiments, provided that the relevant similarity parameters are reproduced, as discussed in Section~\ref{sec:val_flow}.

Since both the configuration and the flow at zero sideslip are symmetric with respect to the vertical plane containing the fuselage axis, only half of the aircraft is simulated and a symmetry condition is imposed on that plane. Aerodynamic coefficients are therefore computed with half of the reference area, so that they correspond to those of the complete aircraft. The reference quantities are summarised in Table~\ref{tab:val_refquantities}.
% TODO: se si calcola C_M, aggiungere: "the pitching moment is evaluated about the moment reference centre of the DPW-IV geometry" + riga in tabella

\begin{table}[htbp]
  \centering
  \caption{Reference quantities of the full-scale CRM.}
  \label{tab:val_refquantities}
  \begin{tabular}{lcc}
    \hline
    Quantity & Value & SI units \\
    \hline
    Reference area, $S_\mathrm{ref}$ & \SI{594720}{in^2} & \SI{383.69}{m^2} \\
    Half-model reference area, $S_\mathrm{ref}/2$ & \SI{297360}{in^2} & \SI{191.845}{m^2} \\
    Mean aerodynamic chord, $c_\mathrm{ref}$ & \SI{275.8}{in} & \SI{7.005}{m} \\
    Wing span, $b$ & \SI{2313.4}{in} & \SI{58.76}{m} \\
    Aspect ratio, $AR$ & \multicolumn{2}{c}{9.0} \\
    \hline
  \end{tabular}
\end{table}

\subsection{Flow conditions}
\label{sec:val_flow}

The simulations reproduce the cruise condition at $M_\infty = 0.85$ and $Re = 5\times10^{6}$, based on $c_\mathrm{ref}$, which is the condition adopted in DPW-IV and at which the CRM has been tested in all the facilities mentioned above. Because the computational geometry is about 37 times larger than the wind tunnel model, matching the experimental Reynolds number at the same Mach number requires a correspondingly lower freestream density. The freestream state has therefore been derived from the two similarity parameters, once the temperature has been fixed.

The freestream static temperature is set to $T_\infty = \SI{310.83}{K}$ (\ang{100}F), the reference temperature of the DPW-IV test case. The dynamic viscosity follows from Sutherland's law,
\begin{equation}
  \mu_\infty = \mu_0 \left( \frac{T_\infty}{T_0} \right)^{3/2} \frac{T_0 + S}{T_\infty + S},
  \label{eq:val_sutherland}
\end{equation}
with $\mu_0 = \SI{1.716e-5}{kg/(m.s)}$, $T_0 = \SI{273.11}{K}$ and $S = \SI{110.56}{K}$, giving $\mu_\infty = \SI{1.897e-5}{kg/(m.s)}$. The freestream velocity is obtained from the Mach number and the speed of sound of the gas, treated as calorically perfect with $\gamma = 1.4$,
\begin{equation}
  U_\infty = M_\infty \sqrt{\gamma R T_\infty},
  \label{eq:val_velocity}
\end{equation}
and the density follows from the definition of the Reynolds number. Combining it with the equation of state yields the freestream pressure,
\begin{equation}
  \rho_\infty = \frac{Re\, \mu_\infty}{U_\infty\, c_\mathrm{ref}}, \qquad p_\infty = \rho_\infty R T_\infty .
  \label{eq:val_pressure}
\end{equation}
The resulting values are listed in Table~\ref{tab:val_flowconditions}. The static pressure, of about \SI{4}{kPa}, has no physical meaning in itself: it is simply the value that, on the full-scale geometry, restores the ratio between inertial and viscous forces of the wind tunnel test.

\begin{table}[htbp]
  \centering
  \caption{Freestream conditions used in the simulations.}
  \label{tab:val_flowconditions}
  \begin{tabular}{lc}
    \hline
    Quantity & Value \\
    \hline
    Mach number, $M_\infty$ & 0.85 \\
    Reynolds number, $Re$ (based on $c_\mathrm{ref}$) & $5\times10^{6}$ \\
    Static temperature, $T_\infty$ & \SI{310.83}{K} \\
    Dynamic viscosity, $\mu_\infty$ & \SI{1.897e-5}{kg/(m.s)} \\
    Velocity, $U_\infty$ & \SI{300.3}{m/s} \\
    Density, $\rho_\infty$ & \SI{0.0451}{kg/m^3} \\
    Static pressure, $p_\infty$ & \SI{4022}{Pa} \\
    \hline
  \end{tabular}
\end{table}

The reference cruise condition is computed at an angle of attack $\alpha = \SI{2.38}{\degree}$, for which the Fluent solution gives a lift coefficient very close to the design value $C_L = 0.5$. The same incidence is used for all the comparisons at cruise conditions, including those between the two turbulence models and between Fluent and STREAmS. The boundary layer is assumed to be fully turbulent: in the experiments transition was fixed by trip dots placed at 10\% of the local chord~\cite{rivers2015}, so that the laminar region is limited to a small portion close to the leading edge.

\subsection{Reference data}
\label{sec:val_refdata}

The experimental data used for the comparison are taken from the ETW test campaign carried out within the ESWIRP project and reported by Rivers et al.~\cite{rivers2015}. Among the measurements available at $M_\infty = 0.85$ and $Re = 5\times10^{6}$, the data point with the lift coefficient closest to the design value was selected (run 153, point 345), corresponding to $\alpha = \SI{2.88}{\degree}$ and $C_L = 0.498$. The comparison covers all the six spanwise stations at which pressure distributions are reported for this point, $\eta = 0.201$, $0.283$, $0.502$, $0.603$, $0.846$ and $0.950$, where $\eta$ is the spanwise coordinate normalised by the semi-span; the measurement uncertainty on $C_p$ is estimated at $\pm 0.0026$~\cite{rivers2015}.

The experimental lift coefficient is reached at an angle of attack about half a degree higher than in the computations. For this reason the comparison is carried out at matched lift coefficient rather than at matched incidence, which is also the criterion followed in~\cite{rivers2015} to compare data from different facilities. The difference in incidence reflects effects that are present in the experiment but not in the simulations, such as the deformation of the wing under load and the interference of the model support, which was not corrected for in the published data~\cite{rivers2015}; this point is discussed further in Section~\ref{sec:val_cruise}.

The numerical reference is provided by Du and Agarwal~\cite{du2017}, who computed the same configuration and flow conditions with Ansys Fluent, using both the Spalart--Allmaras and the SST $k$--$\omega$ turbulence models on a structured grid of about 11 million nodes. Their results include the lift and drag coefficients over the range $\alpha = \SIrange{0}{4}{\degree}$, and are used here mainly for the comparison of the integral coefficients and of their variation with the angle of attack.
